\documentclass[10pt,a4paper]{amsart}
\usepackage[margin=19mm]{geometry}
\usepackage{amsmath,amssymb,mathtools}
\usepackage[scr=rsfs,cal=euler]{mathalfa}
\usepackage{xfrac}
\usepackage[unicode,hidelinks,bookmarks=false]{hyperref}
\newcommand{\ie}{i.e.\ }
\newcommand{\cf}{cf.\ }
\newcommand{\resp}{respectively\ }
\newcommand{\Subsection}{\subsection}
\providecommand{\nobreakdash}{\nobreak\hbox{-}}
\setlength{\emergencystretch}{2em}
\multlinegap0pt
\usepackage{amssymb}
\usepackage{algorithm}\usepackage{algpseudocode}
\usepackage{float}
\usepackage{enumerate}
\usepackage{bm}
\newtheorem{proposition}{Proposition}
\newtheorem{theorem}{Theorem}
\newcommand{\countvar}[4]{z_{#1#2,#3#4}}
\newcommand{\drivevar}[2]{y_{#1#2}}
\newcommand{\entrancex}{p}
\newcommand{\entrancey}{q}
\newcommand{\exitset}{\{(\exitx, \exity)\}}
\newcommand{\exitx}{r}
\newcommand{\exity}{s}
\newcommand{\flowvar}[4]{f_{#1#2,#3#4}}
\newcommand{\flowvardash}[4]{g_{#1#2,#3#4}}
\newcommand{\grid}{G}
\newcommand{\gridadjset}[2]{D_{#1#2}}
\newcommand{\gridarcset}{A}
\newcommand{\gridinvadjset}[2]{D^{-1}_{#1#2}}
\newcommand{\onewaypartition}{W}
\newcommand{\onewaypartitioncomp}{W'}
\newcommand{\onewaypartitionsets}[1]{\mathcal{W}_{#1}}
\newcommand{\validdriveset}{D}
\title{A Branch-and-Cut Algorithm for the Optimal Design of Parking Lots with One-way and Two-way Lanes}
\author{Helen Thomas, Tarun Rambha}
\date{Source: arXiv:2506.09961v3 (CC BY 4.0)}
\begin{document}
\maketitle

\noindent\emph{Source and licence.} A Branch-and-Cut Algorithm for the Optimal Design of Parking Lots with One-way and Two-way Lanes. Version arXiv:2506.09961v3 (CC BY 4.0), distributed by arXiv under the Creative Commons Attribution 4.0 International licence, http://creativecommons.org/licenses/by/4.0/, as recorded in the arXiv licence metadata for this exact version and shown on the abstract page https://arxiv.org/abs/2506.09961v3. Full licence notice: \texttt{licenses/arxiv-2506.09961-CC-BY-4.0.txt}.\par\smallskip

\noindent\textit{Editorial scope.} Counted unit: the article's theorem theorem (source0.tex lines 1628--1630). The article's complete original proof of it is delivered with it (source0.tex lines 1631--1694, one \texttt{proof} environment of 64 lines). No mathematics is written, repaired or re-derived: the statement and the proof are the article's own lines, in the article's own notation, and no retained line is substituted or re-flowed.  Retained uncounted as the source-local context the proof needs: lines 742--745; lines 1163--1166; lines 1169--1172; lines 1175--1180; lines 1184--1188; lines 1503--1506; lines 1525--1532.  Nothing was omitted from any retained span, so the package carries no omission record.  No retained line contains a citation, so the wrapper carries no bibliography and no bibliography entry.  No retained line contains a figure environment, an image-inclusion command or drawing code, so no image is delivered and the package declares no asset dependency.  Editorial formatting only, in addition: an AMS \texttt{amsart} preamble that restates the article's own declarations and macros used by the retained lines, namely the environments \texttt{proposition}, \texttt{theorem} on separate counters, together with the standard packages those lines need.  The retained environments print in this document as Proposition 1, Theorem 1 in order of appearance, whereas the article prints the same items with the numbering of its own full document.  The complete native source of the article as distributed in the arXiv e-print for this exact version is bundled unchanged as \texttt{sources/179702/source0.tex}.
\begin{align} 
    & \countvar{i}{j}{k}{l} \leq \flowvar{i}{j}{k}{l} \leq  M \countvar{i}{j}{k}{l} \qquad && \forall \, ((i,j), (k,l)) \in \gridarcset  \label{eq:count_var_bounds}\\
    & \countvar{i}{j}{k}{l} \leq \flowvardash{i}{j}{k}{l} \leq M \countvar{i}{j}{k}{l} \qquad && \forall \, ((i,j), (k,l)) \in \gridarcset   \label{eq:count_dash_var_bounds}     
\end{align}

\begin{align}
    & \drivevar{m}{n} \leq \sum_{\substack{((i,j),(k,l)) \in \gridarcset: \\ (i,j) \in \onewaypartitioncomp, (k, l) \in \onewaypartition}} \countvar{i}{j}{k}{l} \qquad && \forall \,  \onewaypartition \in \onewaypartitionsets{\entrancex \entrancey}, (m, n) \in \onewaypartitioncomp \label{eq:cutconstraint_entrance}\\
    & \drivevar{m}{n} \leq  \sum_{\substack{((i,j),(k,l)) \in \gridarcset: \\(i,j) \in \onewaypartition, (k, l) \in \onewaypartitioncomp}} \countvar{i}{j}{k}{l} \qquad && \forall \, \onewaypartition \in \onewaypartitionsets{\exitx \exity}, (m, n) \in \onewaypartitioncomp \label{eq:cutconstraint_exit}
\end{align}

\begin{proposition}
\label{prop:1w_aggregate}
Constraints \eqref{eq:cutconstraint_entrance} and \eqref{eq:cutconstraint_exit} are valid for the feasible region of the formulation $\mathcal{F}^{\textsc{flow}}_{\textsc{1W}}$.
\end{proposition}

\begin{align}
\sum_{(m,n) \in \onewaypartitioncomp} \drivevar{m}{n} & = \sum_{\substack{((i,j),(k,l)) \in \gridarcset: \\(i,j) \in \onewaypartitioncomp, (k, l) \in \onewaypartition}} f_{ij, kl} -
\sum_{\substack{((i,j),(k,l)) \in \gridarcset: \\(k, l) \in \onewaypartition, (i,j) \in \onewaypartitioncomp}} f_{kl, ij} \notag \\
& \leq \sum_{\substack{((i,j),(k,l)) \in \gridarcset: \\(i,j) \in \onewaypartitioncomp, (k, l) \in \onewaypartition}} f_{ij, kl} \notag \\
& \leq \sum_{\substack{((i,j),(k,l)) \in \gridarcset: \\(i,j) \in \onewaypartitioncomp, (k, l) \in \onewaypartition}} M \countvar{i}{j}{k}{l} \qquad \text{[using \eqref{eq:count_var_bounds}]} \label{eq:aggregate_vi_1w_a}
\end{align}

\begin{align}
\sum_{(m,n) \in \onewaypartitioncomp} \drivevar{m}{n} & = \sum_{\substack{((i,j),(k,l)) \in \gridarcset: \\(i, j) \in \onewaypartition, (k,l) \in \onewaypartitioncomp}} g_{ij, kl} - \sum_{\substack{((i,j),(k,l)) \in \gridarcset: \\(k,l) \in \onewaypartitioncomp, (i, j) \in \onewaypartition}} g_{kl, ij} \notag \\
& \leq \sum_{\substack{((i,j),(k,l)) \in \gridarcset: \\(i, j) \in \onewaypartition, (k,l) \in \onewaypartitioncomp}} g_{ij, kl}  \notag \\
& \leq \sum_{\substack{((i,j),(k,l)) \in \gridarcset: \\(i, j) \in \onewaypartition, (k,l) \in \onewaypartitioncomp}} M \countvar{i}{j}{k}{l} \qquad \text{[using \eqref{eq:count_dash_var_bounds}]} \label{eq:aggregate_vi_1w_b}
\end{align}

\begin{theorem}
\label{theorem:two_way}
Suppose $(x, y)$ denotes an integer feasible parking and driving-field assignment. The formulations $\mathcal{F}_{\textsc{2W}}^{\textsc{flow}}$ and $\mathcal{F}_{\textsc{2W}}^{\textsc{CCC}}$ induce the same set of integer feasible values for $(x,y)$.
\end{theorem} 

\begin{align}
    \text{s.t. } & \alpha_{ij} - \alpha_{kl} + \beta_{ij, kl} \geq 0 && \forall \, ((i,j), (k,l)) \in \gridarcset, (i, j), (k, l) \neq (p, q)  \label{eq:sp_dual_1} \\
     & - \alpha_{kl} + \beta_{pq, kl} \geq 0 && \forall \,  (k, l) \in \gridadjset{p}{q}  \label{eq:sp_dual_2} \\
    & \alpha_{pq} + \beta_{ij, pq} \geq 0 && \forall \,  (i, j) \in \gridinvadjset{p}{q}  \label{eq:sp_dual_3} \\
    & \alpha_{ij} \leq 0 && \forall \, (i,j) \in \validdriveset \setminus \{(\entrancex, \entrancey)\} \label{eq:sp_dual_4} \\
    & \alpha_{pq} \geq 0 &&  \label{eq:sp_dual_4.5} \\
    & \beta_{ij, kl} \geq 0 && \forall \, ((i,j), (k,l)) \in \gridarcset \label{eq:sp_dual_5} 
\end{align}

\begin{theorem}
Suppose $(x, y, z)$ denotes an integer feasible parking, driving-field, and direction assignment. The formulations $\mathcal{F}_{\textsc{1W}}^{\textsc{flow}}$ and $\mathcal{F}_{\textsc{1W}}^{\textsc{CCC}}$ induce the same set of integer feasible values for $(x,y,z)$.
\end{theorem}

\begin{proof}
We again write the Benders reformulation of the problem $\mathcal{F}_{\textsc{1W}}^{\textsc{flow}}$ using two subproblems, since the constraints with $f$ and $g$ variables have a block diagonal structure for a fixed $(x, y, z)$. 
\begin{align}
    \max & \, \, \, \, 0  && \notag \\
      \text{s.t. }  & \sum_{(k,l) \in \gridadjset{i}{j}} \flowvar{i}{j}{k}{l} - \sum_{(k,l) \in \gridinvadjset{i}{j}} \flowvar{k}{l}{i}{j} \geq \drivevar{i}{j} \qquad && \forall \, (i,j) \in \validdriveset \setminus \{(\entrancex, \entrancey)\} \label{eq:floworigin_1w_sp}\\
     & \sum_{(i,j) \in \gridinvadjset{\entrancex}{\entrancey}} \flowvar{i}{j}{\entrancex}{\entrancey} \leq \sum_{(i,j) \in \validdriveset} \drivevar{i}{j} -1 \qquad &&  \label{eq:flowdest_1w_sp}\\
    & \flowvar{i}{j}{k}{l} \leq M \countvar{i}{j}{k}{l} \qquad && \forall \, ((i,j), (k,l)) \in \gridarcset \label{eq:flowhead_1w_sp} \\
    & \flowvar{i}{j}{k}{l} \geq 0 \qquad && \forall \, ((i,j), (k,l)) \in \gridarcset \label{eq:flowvar_1w_nonneg}
\end{align}
Let $\alpha$ be the dual variables for \eqref{eq:floworigin_1w_sp}--\eqref{eq:flowdest_1w_sp}, and $\beta$ denote the duals for \eqref{eq:flowvar_1w_nonneg}. The dual of the first subproblem is
\begin{align}
    \min & \sum_{(i,j) \in \validdriveset \setminus \{(\entrancex, \entrancey)\}} \drivevar{i}{j} \alpha_{ij} + \Bigg(\sum_{(i,j) \in \validdriveset} \drivevar{i}{j} -1 \Bigg) \alpha_{pq} + \sum_{((i,j), (k,l)) \in \gridarcset}  M \countvar{i}{j}{k}{l} \beta_{ij, kl}  \\
    \text{s.t. } & \eqref{eq:sp_dual_1} \text{--} \eqref{eq:sp_dual_5}
\end{align}
Note that the feasible region is identical to that of the dual subproblem in Theorem \ref{theorem:two_way}. Hence, the non-trivial extreme directions in $E$ yield feasibility cuts of the form
\begin{align}
\sum_{(i,j) \in S} \drivevar{i}{j} \leq  \sum_{\substack{
  ((i,j),(k,l)) \in \gridarcset:\\
  (i,j) \in S,\,
  (k,l) \in D \setminus S
}} M \countvar{i}{j}{k}{l}  \label{eq:simplified_cut_1w_a}
\end{align}
Similarly, the second subproblem can be written as 
\begin{align}
    \max & \, \, \, \, 0  && \notag \\
     \text{s.t. } & \sum_{(k,l) \in \gridinvadjset{i}{j}} \flowvardash{k}{l}{i}{j} - \sum_{(k,l) \in \gridadjset{i}{j}} \flowvardash{i}{j}{k}{l} \geq \drivevar{i}{j}  \qquad && \forall \, (i,j) \in \validdriveset \setminus \exitset \label{eq:flow_dash_origin_singlelane_sp}\\
    & \sum_{(i,j) \in \gridadjset{\exitx}{\exity}} \flowvardash{\exitx}{\exity}{i}{j} \leq \sum_{(i,j) \in \validdriveset} \drivevar{i}{j}  - 1 \qquad && \label{eq:flowdashdest_singlelane_sp}\\
    & \flowvardash{i}{j}{k}{l} \leq M \countvar{i}{j}{k}{l}  \qquad && \forall \, ((i,j), (k,l)) \in \gridarcset \label{eq:flow_dash_tail_singlelane_sp}\\
    & \flowvardash{i}{j}{k}{l} \geq 0 \qquad && \forall \, ((i,j), (k,l)) \in \gridarcset \label{eq:g_non_negative_sp}
\end{align}
Let $\pi$ denote the dual variables for \eqref{eq:flow_dash_origin_singlelane_sp}--\eqref{eq:flowdashdest_singlelane_sp} and $\xi$ represent the dual variables for \eqref{eq:flow_dash_tail_singlelane_sp}. 
The dual subproblem for the above formulation is given by 
\begin{align}
    \min &\sum_{(i,j) \in \validdriveset \setminus \exitset} \drivevar{i}{j}\pi_{ij} + \Bigg(\sum_{(i,j) \in \validdriveset} \drivevar{i}{j}  - 1 \Bigg) \pi_{rs} + \sum_{((i,j), (k,l)) \in \gridarcset}  M \countvar{i}{j}{k}{l} \xi_{ij, kl} \label{eq:sp_dual_1w}
\end{align}
\begin{align}
     \text{s.t. } & - \pi_{ij} + \pi_{kl} + \xi_{ij, kl} \geq 0 && \forall \, ((i,j), (k,l)) \in \gridarcset, (i, j), (k, l) \neq (r, s)  \label{eq:sp_dual_1w_6} \\
     & - \pi_{ij} + \xi_{ij, rs} \geq 0 && \forall \,  (i,j) \in \gridinvadjset{r}{s}  \label{eq:sp_dual_1w_7} \\
    & \pi_{rs} + \xi_{rs, kl} \geq 0 && \forall \,  (k,l) \in \gridadjset{r}{s}  \label{eq:sp_dual_1w_8} \\
    & \pi_{ij} \leq 0 && \forall \, (i,j) \in \validdriveset \setminus \{(r, s)\} \label{eq:sp_dual_1w_9} \\
    & \pi_{rs} \geq 0 && \label{eq:sp_dual_1w_9.5} \\
     & \xi_{ij, kl} \geq 0 && \forall \, ((i,j), (k,l)) \in \gridarcset \label{eq:sp_dual_1w_10} 
\end{align}
The extreme directions of the pointed cone induced by \eqref{eq:sp_dual_1w_6}--\eqref{eq:sp_dual_1w_10} are given below. The justification is identical to that in Theorem \ref{theorem:two_way} and is therefore omitted.
\begin{enumerate}
\item A singleton $F_1$ consisting of a vector with $\pi_{rs} = 1$, and all remaining variables equal to zero.
\item A set of directions $F_2$, one for each arc $((i,j),(k,l)) \in \gridarcset$, with $\xi_{ij,kl} = 1$ and all other entries equal to zero.
\item A collection of vectors $F_3$, where each vector is defined by a subset 
$S \subseteq \validdriveset \setminus \{(r,s)\}$ such that the induced subgraph 
$\grid(S)$ is connected, and has components 
$\pi_{ij} = -1$ for all $(i,j) \in S$, 
$\xi_{ij,kl} = 1$ for all $((i,j),(k,l)) \in \gridarcset$ with $(k,l) \in S$ and $(i,j) \in \validdriveset \setminus S$, 
and all remaining components equal to zero.
\end{enumerate}
The extreme directions from $F_1 \cup F_2$ produce trivial feasibility cuts. The cuts from those in $F_3$ can be written as
\begin{align}
\sum_{(i,j) \in S} \drivevar{i}{j} \leq  \sum_{\substack{
  ((i,j),(k,l)) \in \gridarcset:\\
  (i,j) \in D \setminus S,\,
  (k,l) \in S
}} M \countvar{i}{j}{k}{l}  \label{eq:simplified_cut_1w_b}
\end{align}
Constraints \eqref{eq:simplified_cut_1w_a} and \eqref{eq:simplified_cut_1w_b} are the same as the valid inequalities \eqref{eq:aggregate_vi_1w_a} and \eqref{eq:aggregate_vi_1w_b}, respectively. Together with Proposition \ref{prop:1w_aggregate}, this implies that the integer feasible regions of $\mathcal{F}_{\textsc{1W}}^{\textsc{flow}}$ and $\mathcal{F}_{\textsc{1W}}^{\textsc{CCC}}$ coincide.
\end{proof}

\end{document}
